\begin{name}
	{\tenchude}
	{TOÁN 10}
	{LỚP TOÁN THẦY PHÁT}
	{Life always offers you a second chance. It’s called tomorrow}
\end{name}
\Opensolutionfile{ansbook}[ans/ansbookDe5]
\TN
\Opensolutionfile{ans}[ans/ansDe5-TN1]
\begin{ex}%[0D6N3-3]
Chỉ số IQ của một nhóm học sinh là
\begin{center}
\begin{tabular}{|l|l|l|l|l|l|l|l|l|l|}
\hline
$60$ & $78$ & $80$ & $64$ & $70$ & $76$ & $80$ & $74$ & $86$ & $90$\\
\hline
\end{tabular}
\end{center}
Tứ phân vị thứ ba của mẫu số liệu là
\choice
{\True $Q_{3}=80$}
{$Q_{3}=70$}
{$Q_{3}=83$}
{$Q_{3}=75$}
\loigiai{
Sắp xếp các giá trị này theo thứ tự không giảm
\begin{center}
\begin{tabular}{|l|l|l|l|l|l|l|l|l|l|}
\hline
$60$ & $64$ & $70$ & $74$ & $76$ & $78$ & $80$ & $80$ & $86$ & $90$ \\
\hline
\end{tabular}
\end{center}
Vì $n=10$ là số chẵn nên $Q_{2}$ là số trung bình cộng của hai giá trị chính giữa $Q_{2}=\dfrac{(76+78)}{2}=77$.
Ta tìm $Q_{3}$ là trung vị của nửa số liệu bên phải $Q_{2}$
\begin{center}
\begin{tabular}{|l|l|l|l|l|}
\hline
$78$ & $80$ & $80$ & $86$ & $90$ \\
\hline
\end{tabular}
\end{center}
và tìm được $Q_{3}=80$.
}
\end{ex}

\begin{ex}%[0D6N4-2]
Hãy tìm khoảng biến thiên của mẫu số liệu thông kê sau
\[22\qquad 24\qquad 33\qquad 17\qquad 11\qquad 4\qquad 18\qquad 87\qquad 72\qquad 30\]
\choice
{$11$}
{$33$}
{$87$}
{\True 83}
\loigiai{
Khoảng biến thiên $R=x_{\max}-x_{\min}=87-4=83$.
}
\end{ex}

\begin{ex}%[0D8H1-3]
Bình $A$ chứa $3$ quả cầu xanh, $4$ quả cầu đỏ và $5$ quả cầu trắng. Bình $B$ chứa $4$ quả cầu xanh, $3$ quả cầu đỏ và $6$ quả cầu trắng. Bình $C$ chứa $5$ quả cầu xanh, $5$ quả cầu đỏ và $2$ quả cầu trắng. Từ mỗi bình lấy một quả cầu. Có bao nhiêu cách lấy để cuối cùng được $3$ quả có màu giống nhau.
\choice
{\True $180$}
{$150$}
{$120$}
{$60$}
\loigiai{
\begin{itemize}
\item Số cách lấy được ba quả cầu màu xanh là $3\cdot 4\cdot 5=60$.
\item Số cách lấy được ba quả cầu màu đỏ là $4\cdot 3\cdot 5=60$.
\item Số cách lấy được ba quả cầu màu trắng là $5\cdot 6\cdot 2=60$.
\end{itemize}
Vậy số cách lấy được ba quả cầu cùng màu là $60+60+60=180$.
}
\end{ex}

\begin{ex}%[0D8H2-3]
Có bao nhiêu số tự nhiên có sáu chữ số khác nhau từng đôi một, trong đó chữ số $ 5$ đứng liền giữa hai chữ số $ 1$ và $ 4$?
\choice
{$249$}
{\True $1500$}
{$3204$}
{$2942$}
\loigiai{
Chữ số $ 5$ đứng liền giữa hai chữ số $ 1$ và $ 4$ nên ta có thể có $ 154$ hoặc $ 451$.\\
Gọi số cần tìm là $\overline{abc}$, sau đó ta chèn thêm $ 154$ hoặc $451$ để có được số gồm 6 chữ số cần tìm.
\begin{itemize}
\item \textbf{Trường hợp 1:} $ a\ne 0$, số cách chọn $ a$ là $ 6$, số cách chọn $ b$ và $ c$ là $ \mathrm{A}_6^2$, sau đó chèn $ 154$ hoặc $ 451$ vào $ 4$ vị trí còn lại nên có $ 6\cdot\mathrm{A}_6^2\cdot4\cdot2$ cách.
\item \textbf{Trường hợp 2:} $ a=0$, số cách chọn $ a$ là 1, số cách chọn $ b$ và $ c$ là $ \mathrm{A}_6^2$, sau đó chèn $ 154$ hoặc $ 451$ vào vị trí trước $ a$ có duy nhất 1 cách nên có $ \mathrm{A}_6^2\cdot2$ cách.
\end{itemize}
Vậy có $ 6\cdot\mathrm{A}_6^2\cdot4\cdot2+
\mathrm{A}_6^2\cdot2=1500$.}
\end{ex}

\begin{ex}%[HSA - đợt 1, Thái Bảo]%[0D8H2-4]
Một tổ chăm sóc khách hàng của một trung tâm điện tử gồm $12$ nhân viên. Số cách phân công $3$ nhân viên đi đến ba địa điểm khác nhau để chăm sóc khách hàng là
\choice
{\True $1320$}
{$1230$}
{$220$}
{$1728$}
\loigiai{
Số cách xếp $3$ nhân viên từ $12$ nhân viên vào $3$ vị trí khác nhau là $\mathrm{A}_{12}^3=1320$ cách.
}
\end{ex}

\begin{ex}%[0D8H2-6]
Cho $10$ điểm phân biệt và không có ba điểm nào thẳng hàng. Hỏi có bao nhiêu véc-tơ khác $\overrightarrow{0}$ có điểm đầu và điểm cuối là hai trong số $10$ điểm đã cho?
\choice
{\True $90$}
{$45$}
{$10$}
{$20$}
\loigiai{
Số véc-tơ khác $\overrightarrow{0}$ có điểm đầu và điểm cuối là hai trong số $10$ điểm đã cho là $\mathrm{A}_{10}^2=90$.
}
\end{ex}

\begin{ex}%[0D8H3-2]
Khai triển $(x+1)^4$ ta được
\choice
{\True $x^4+4x^3+6x^2+4x+1$}
{$x^4-4x^3+6x^2-4x+1$}
{$x^4+6x^3+8x^2+6x+1$}
{$x^4+4x^3+8x^2+4x+1$}
\loigiai{
Ta có $(x+1)^4=\mathrm{C}_4^0\,x^4+\mathrm{C}_4^1\,x^3+\mathrm{C}_4^2\,x^2+\mathrm{C}_4^3\,x+\mathrm{C}_4^4=x^4+4x^3+6x^2+4x+1$.
}
\end{ex}

\begin{ex}%[HK1, THPT Lê Quảng Chí, 2023]%[TVN-001, 10-11EX-HK1-2324]%[0H9H1-2]
Trong hệ tọa độ $Oxy$, cho ba điểm $A(2;1)$, $B(0;-3)$, $C(3;1)$. Tìm tọa độ điểm $D$ để $ABCD$ là hình bình hành.
\choice
{\True $(5;5)$}
{$(5;-2)$}
{$(5;-4)$}
{$(-1;-4)$}
\loigiai{
Tọa độ điểm $D$ cần tìm là
\[\heva{&x_D=x_A+x_C-x_B=2+3-0=5\\&y_D=y_A+y_C-y_B=1+1-(-3)=5.}\]}
\end{ex}

\begin{ex}%[0H9H3-2]%[10-11-12EX-HK1-2425]%[Đỗ Đường Hiếu]
Trong mặt phẳng $Oxy$, cho điểm $M(2;1)$. Gọi $A$, $B$ là hình chiếu của $M$ lên $Ox$, $Oy$. Phương trình đường thẳng $AB$ là
\choice
{$x+y-3=0$}
{$2x+y+2=0$}
{$2x+y-2=0$}
{\True $x+2y-2=0$}
\loigiai{
Ta có $A(2;0)$, $B(0;1)$. Do đó, phương trình đường thẳng $AB$ là
\[\dfrac{x}{2}+\dfrac{y}{1}=1\Leftrightarrow x+2y-2=0.\]
}
\end{ex}

\begin{ex}%[0H9H4-2]%[Dự án đề kiểm tra Toán khối 10 HKI NH23-24-Dot 15- Hồ Văn Trung]%[THPT Lê Qúy Đôn- TPHCM]
Trong mặt phẳng $Oxy$, phương trình đường tròn có tâm $I\left(3;1\right)$ và đi qua điểm $M\left(2;-1\right)$ là
\choice
{$\left(x+3\right)^2+\left(y+1\right)^2=5$}
{$\left(x-3\right)^2+\left(y-1\right)^2=\sqrt{5}$}
{$\left(x+3\right)^2+\left(y+1\right)^2=\sqrt{5}$}
{\True $\left(x-3\right)^2+\left(y-1\right)^2=5$}
\loigiai{Ta có $IM=\sqrt{(2-3)^2+(-1-1)^2}=\sqrt{5}$.\\
Đường tròn tâm $I(3;1)$ và có $R=IM=\sqrt{5}$ là
$(C) \colon (x-3)^2+(y-1)^2=5$.
}
\end{ex}

\begin{ex}%[0H9H5-2]
Trong mặt phẳng $Oxy$, phương trình chính tắc của Elip có độ dài  trục lớn  bằng $6$ và tỉ số giữa tiêu cự với độ dài trục lớn bằng $\dfrac{1}{3}$ là
\choice
{$\dfrac{x^2}{9}+\dfrac{y^2}{3}=1$}
{$\dfrac{x^2}{9}+\dfrac{y^2}{5}=1$}
{\True $\dfrac{x^2}{9}+\dfrac{y^2}{8}=1$}
{$\dfrac{x^2}{6}+\dfrac{y^2}{5}=1$}
\loigiai{ Độ dài trục lớn $2a=6\Rightarrow a=3$.\\
Tỉ số giữa tiêu cự và độ dài trục lớn  $\dfrac{2c}{2a}=\dfrac{1}{3}\Rightarrow c =\dfrac{1}{3} \cdot a=\dfrac{1}{3}\cdot 3=1$.\\
Ta có $c^2=a^2-b^2\Leftrightarrow 1^2=3^2-b^2\Leftrightarrow b^2=3^2-1^2=8$. \\
Phương trình chính tắc của Elip $\dfrac{x^2}{9}+\dfrac{y^2}{8}=1$.
}
\end{ex}

\begin{ex}%[0H9H5-4]
Tổng các giá trị $m$ để hypebol $(H)\colon \dfrac{x^2}{4}-\dfrac{y^2}{m^2-2m-7}=1$ có tiêu điểm trung với tiêu điểm của elip $(E)\colon\dfrac{x^2}{16}+\dfrac{y^2}{9}=1$ là
\choice
{$4$}
{\True $2$}
{$5$}
{$6$}
\loigiai{
Tiêu điểm của elip $(E)\colon \dfrac{x^2}{16}+\dfrac{y^2}{9}=1$ là
\[F_{1;2}=\left(\pm\sqrt{16-9};0\right)=\left(\pm\sqrt{7};0\right).\]
Tiêu điểm của hypebol $(H)\colon \dfrac{x^2}{4}-\dfrac{y^2}{m^2-2m-7}=1$ là \[F'_{1;2}=\left(\pm\sqrt{a^2+b^2};0\right)=\left(\pm\sqrt{4+m^2-2m-7};0\right).\]
\begin{eqnarray*}
&\Rightarrow& \left(\pm\sqrt{4+m^2-2m-7};0\right)=\left(\pm\sqrt{7};0\right)\\
&\Rightarrow& 4+m^2-2m-7=7\\
&\Rightarrow& m^2-2m-10=0\\
&\Rightarrow& \hoac{&m=1+\sqrt{11}\\ &m=1-\sqrt{11}.}
\end{eqnarray*}
Vậy tổng của giá trị $m$ để thỏa yêu cầu bài toán là $2$.
}
\end{ex}

\TNTF
\begin{ex}%[0D8H2-4]
Lớp 10A có $30$ học sinh gồm $13$ bạn nữ và $17$ bạn nam, trong đó bạn nam tên Khang làm lớp trưởng.
\choiceTF
{Chọn ra hai bạn gồm một nam và một nữ tham gia vào Đội cờ đỏ. Số cách chọn là $220$ cách}
{\True Chọn ra ba bạn trực nhật lớp, trong đó phân công một bạn quét lớp, một bạn quét sân và một bạn lau bảng. Số cách chọn là $24360$}
{\True Chọn ra ba bạn tham gia hoạt động thiện nguyện, trong đó phải có lớp trưởng và có ít nhất một nữ. Số cách chọn là $286$}
{Sắp xếp học sinh để chụp ảnh kỉ yếu trong đó có $14$ bạn đứng hàng trước và $16$ bạn đứng hàng sau. Số cách sắp xếp là $14!\cdot 16!$}
\loigiai
{
\begin{itemchoice}
\itemch Số cách chọn ra hai bạn gồm một nam và một nữ là $13\cdot 17=221$ cách.
\itemch Số cách chọn $3$ bạn từ $30$ bạn và có sắp xếp theo công việc là $\mathrm{A}_{30}^3=24360$.
\itemch Chọn ra ba bạn tham gia hoạt động thiện nguyện, trong đó phải có lớp trưởng và có ít nhất một nữ.
\begin{enumerate}[TH 1.]
\item Chọn $1$ bạn nữ và bạn Khang cùng $1$ bạn nam, số cách chọn $13\cdot 16=208$.
\item Chọn $2$ bạn nữ và bạn Khang lớp trưởng, số cách chọn $\mathrm{C}_{13}^{2}=78$.
\end{enumerate}
Vậy số cách chọn là $208+78=286$.
\itemch Số cách sắp xếp là $\mathrm{A}_{30}^{14}\cdot 16!$.
\end{itemchoice}
}
\end{ex}

\begin{ex}%[Dự án 18 - TeamTeXHoa - Phạm Quốc Toàn]%[0H9H4-5]
Trong hệ trục tọa độ $Oxy$, cho đường tròn $\left(C \right)$ tâm $I \left(1;2 \right)$ và cắt đường thẳng \\
$\Delta \colon 3x+4y-6=0$ tại hai điểm $A$, $B$ sao cho $S_{IAB}=4$. Trong mỗi ý a), b), c), d) sau, hãy chọn đúng hoặc sai.
\choiceTF
{\True Khoảng cách từ tâm $I$ đến đường thẳng $\Delta$ bằng $1$}
{Bán kính đường tròn $\left(C \right)$ nhỏ hơn $4$}
{Phương trình đường tròn $\left(C \right) \colon x^2+y^2-2x-4y+12=0$}
{Điểm $O$ nằm trên đường tròn $\left(C \right)$}
\loigiai{
\begin{center}
\begin{tikzpicture}[>=stealth, x=1.0 cm, y=1.0 cm]
% \draw[->] (-4,0) --(0,0) node[below left]{$O$} -- (4,0) node[below]{$x$};
% \draw[->](0,-4)--(0, 4) node[right]{$y$};
\draw(0,0) circle (3); %(0,0) là tọa độ tâm, 1 là bán kính
\draw(-4,2) -- (4,2) (0, 0) -- (0, 2) (0, 0) -- (-2.24, 2) (0, 0) -- (2.24, 2);
\draw (0,0) circle (1pt) node[below]{$I$};
\draw (0,2) circle (1pt) node[above]{$H$};
\draw (3.5,2) circle node[below]{$\Delta$};
\draw (-2.24,2) circle (1pt) node[above left]{$A$};
\draw (2.24,2) circle (1pt) node[above right]{$B$};
\end{tikzpicture}
\end{center}
\begin{itemchoice}
\itemch Gọi $H$ là trung điểm của đoạn thẳng $AB \Rightarrow IH \perp AB$. Ta có $IH=d_{\left(I, \Delta\right)}=\dfrac{\left| 3+8-6\right|}{\sqrt{3^2+4^2}}=1$.
\itemch $IH \perp AB \Rightarrow S_{IAB}=\dfrac{1}{2} AB\cdot IH \Leftrightarrow 4=\dfrac{1}{2}AB\cdot 1\Leftrightarrow AB=8$ \\
$\Rightarrow AH=4$ và $R=\sqrt{AH^2+IH^2}=\sqrt{17}$ $\Rightarrow R>4$.
\itemch Phương trình đường tròn $\left(C \right)$ tâm $I \left(1;2 \right)$, bán kính $R=\sqrt{17}$ là $\left(C \right) \colon \left(x-1 \right)^2+ \left(y-2 \right)^2=17$ hay $\left(C \right) \colon x^2+y^2-2x-4y-12=0$.
\itemch Ta có $\overrightarrow{IO}= \left(-1;-2 \right) \Rightarrow IO=\sqrt{\left(-1 \right)^2+ \left(-2 \right)^2}=\sqrt{5}<R=\sqrt{17}$. Suy ra $O$ nằm trong $\left(C \right)$.
\end{itemchoice}
}
\end{ex}

\TNSA



\begin{ex}%[0D8H2-5]%[Dự án Toán thực tế]%[Vũ Hồng Toàn]
Mỗi máy tính tham gia vào mạng phải có một địa chỉ duy nhất, gọi là địa chỉ IP, nhằm định danh máy tính đó trên Internet. Xét tập hợp $A$ gồm các địa chỉ IP có dạng $192$.$168$$ abc$.$deg$, trong đó $a$, $d$ là các chữ số khác nhau được chọn ra từ các chữ số $1$, $2$ còn $b$, $c$, $e$, $g$ là các chữ số đôi một khác nhau được chọn ra từ các chữ số $0,1,2,3,4,5$. Hỏi tập hợp $A$ có bao nhiêu phần tử?
\shortans{$1\,800$}
\loigiai{
$a$, $d$ là các chữ số khác nhau được chọn ra từ các chữ số $1$, $2$ nên có $2$ cách chọn.\\
$b$ và $c$ là các chữ số đôi một khác nhau được chọn ra từ các chữ số $0,1,2,3,4,5$ nên có $\mathrm{A}_6^2$ cách chọn.\\
$e$, $g$ là các chữ số đôi một khác nhau được chọn ra từ các chữ số $0,1,2,3,4,5$ nên có $\mathrm{A}_6^2$ cách chọn.\\
Vậy tập hợp $A$ có số phần tử là $2\cdot \mathrm{A}_6^2\cdot 1\cdot \mathrm{A}_6^2=1\,800$ (phần tử).
}
\end{ex}

\begin{ex}%[0D8H3-4]%[Dự án đề kiểm tra Toán 10, 11 HKII NH23-24-BCTuan]%[THPT LTV-Hà Nội]
Tìm hệ số của số hạng chứa $x^2y^3$ trong khai triển $(2x+y)^5$.
\shortans[3]{$40$}
\loigiai{
Ta có $(2x+y)^5=\displaystyle\sum\limits_{k=0}^5 \mathrm{C}_5^k2^{5-k}x^{5-k}y^k$.\\
Số hạng chứa $x^2y^3$ tương ứng với $k=3$.\\
Hệ số của số hạng là $\mathrm{C}_5^2 2^2=40$.
}
\end{ex}

\begin{ex}%[0H9V4-3]%[Dự án đề kiểm tra Toán 11 GHKII NH23-24- Tổng Nguyễn]%[THPT Đống Đa - Hà Nội]
Trong mặt phẳng $Oxy$, đường tròn $(C)\colon (x-a)^2+(y-b)^2=25$ tiếp xúc với đường thẳng $d\colon 3x+4y+1=0$ tại điểm $A(1;-1)$. Tính  $\dfrac{a}{b}$ (biết $a<0$).
\shortans[]{$0{,}4$}
\loigiai{Đường thẳng $d$ có vectơ pháp tuyến $\overrightarrow{n}=(3;4)$.\\
Đường thẳng $\Delta$ đi qua $A(1;-1)$ và vuông góc với $d$ nên có vectơ chỉ phương $\overrightarrow{u}=(3;4)$. Phương trình của $\Delta$ là
\[\heva{&x=1+3t\\&y=-1+4t} \quad t \in \mathbb{R}.\]
Đường tròn $(C)$ có bán kính $R=5$,  tiếp xúc với $d$ tại $A(1;-1)$ nên tâm $I$ thuộc $\Delta$. Suy ra $I(1+3t; -1+4t)$.\\
Ta có \[ IA=R\Leftrightarrow \sqrt{(-3t)^2+(-4t)^2}=5\Leftrightarrow 25t^2=25 \Leftrightarrow  \hoac{&t=1\Rightarrow I_1(4;3) \ \text{(loại)}\\&t=-1\Rightarrow I_2(-2;-5)\ \text{(nhận)}.}\]
Vậy   $\dfrac{a}{b}=\dfrac{-2}{-5}=\dfrac{2}{5}=0{,}4$.
}
\end{ex}
\begin{ex}%[0H9V5-0]%[Làm đề GHK, HK, 4 phần - 2025]%[Huỳnh Xuân Tín]
Để chụp toàn cảnh, ta có thể sử dụng một gương hypebol. Máy ảnh được hướng về phía đỉnh của gương và tâm quang học của máy ảnh được đặt tại một tiêu điểm của gương (xem hình). Tìm khoảng cách từ tâm quang của máy ảnh đến đỉnh của gương, biết rằng phương trình mặt cắt của gương là $\dfrac{x^2}{25}-\dfrac{y^2}{16}=1$.
\shortans{$11{,}2$}
\begin{center}
\begin{tikzpicture}[scale=0.6, font=\footnotesize, line join=round, line cap=round, >=stealth]
\def\xmin{-10}\def\xmax{10}\def\ymin{-5}\def\ymax{5}\pgfmathsetmacro\c{2*sqrt(5)}
\draw[->] (\xmin-0.2,0)--(\xmax+0.2,0) node[below] {\footnotesize $x$};
\draw[->] (0,\ymin-0.2)--(0,\ymax+0.2) node[right] {$y$};
\clip (\xmin,\ymin) rectangle (\xmax,\ymax);
\path
(0,0)coordinate (O)
(-2,0) coordinate (A_1)
(2,0) coordinate (A_2)
(-\c,0) coordinate (F_1)
(\c,0) coordinate (F_2)
(-6.3,3) coordinate (M)
(-6.3,-3) coordinate (N)
;
\draw[blue] (M)--(N);
\fill[blue!50] (F_2)circle (0.3cm);
\fill[blue!10] (M)--(N)--plot ({-2*sqrt(\x*\x+1)},{\x})--(M);
\draw[smooth,variable=\x,blue] plot ({2*sqrt(\x*\x+1)},{\x});
\draw[smooth,variable=\x,blue] plot ({-2*sqrt(\x*\x+1)},{\x});
\foreach \x/\g in {F_1/-90,F_2/-90}
\fill (\x) circle (4.5pt) ;
\foreach \x/\g in {F_1/-90,F_2/-90,A_1/-40,A_2/240,O/-120}
\fill (\x) circle (1.5pt) ($(\x)+(\g:7mm)$) node{$\x$};
\draw (M)--($(M)+(30:1)$)node[above]{gương hypebol};
\draw (F_2)--($(F_2)+(130:3)$)node[above]{tâm máy ảnh};
\draw[dashed](A_1)--(F_1);
\end{tikzpicture}
\end{center}
\loigiai{ Gọi $(H)\colon \dfrac{x^2}{25}-\dfrac{y^2}{16}=1\Rightarrow \heva{&a^2=25\\&b^2=16}\Rightarrow \heva{&a=5\\&b=4}\Rightarrow c=\sqrt{a^2+b^2}=\sqrt{39}$.\\
Tiêu điểm của gương là $F_1(-\sqrt{39};0)$ và $F_2(\sqrt{39};0)$.\\
Đỉnh của gương là $A_1(-5;0)$.\\
Vậy khoảng cách từ tâm của máy ảnh tới đỉnh của gương là\\ $F_2A_1=OA_1+OF_2=5+\sqrt{39}\approx 11{,}2$.
}
\end{ex}
\TL
\begin{ex}%[10-11-GK2-2425]%[Đỗ Minh Phúc]%[0D8H3-2]
Khai triển biểu thức $\left(2 x-\dfrac{3}{2}\right)^5$ bằng nhị thức Newton.
\loigiai
{
Ta có
\begin{eqnarray*}
\left(2x-\dfrac{3}{2}\right)^5&=&\mathrm{C}_{5}^{0}(2x)^5+\mathrm{C}_{5}^{1}(2x)^4\left(-\dfrac{3}{2}\right)+\mathrm{C}_{5}^{2}(2x)^3\left(-\dfrac{3}{2}\right)^2+\mathrm{C}_{5}^{3}(2x)^2\left(-\dfrac{3}{2}\right)^3+\mathrm{C}_{5}^{4}(2x)\left(-\dfrac{3}{2}\right)^4+\mathrm{C}_{5}^{5}\left(-\dfrac{3}{2}\right)^5\\
&=&32x^5-120x^4+180x^3-135x^2+\dfrac{405}{8}x-\dfrac{243}{32}.
\end{eqnarray*}

}
\end{ex}

\begin{ex}%[0D8V2-4]
Một đội thanh niên tình nguyện có $15$ người gồm $12$ nam và $3$ nữ. Hỏi có bao nhiêu cách phân công đội thanh niên tình nguyện đó về giúp đỡ $3$ tỉnh miền núi, sao cho mỗi tỉnh có $4$ nam và $1$ nữ?
\loigiai{
Ta có $\mathrm{C}_3^1\cdot\mathrm{C}_{12}^4$ cách phân công các thanh niên về tỉnh thứ nhất.\\
Với mỗi cách này thì có $\mathrm{C}_2^1 \cdot \mathrm{C}_8^4$ cách phân công số thanh niên còn lại về tỉnh thứ hai.\\
Với mỗi cách phân công trên thì có $\mathrm{C}_1^1\cdot\mathrm{C}_{4}^4$ cách phân công số thanh niên còn lại về tỉnh thứ ba.\\
Do đó, ta có $\mathrm{C}_3^1\cdot\mathrm{C}_{12}^4\cdot\mathrm{C}_2^1\cdot\mathrm{C}_{8}^4\cdot\mathrm{C}_1^1\cdot\mathrm{C}_{4}^4=207900$ cách phân công thỏa mãn đề bài.\\
Gọi số cần lập có dạng $abc$.\\
Chọn $c$ với $c \in\{4;6;8\}$ có $3$ cách.\\
Chọn $\overline{ab}$ có $\mathrm{A}_5^2$ cách.\\
Theo quy tắc nhân, ta có $3 \cdot \mathrm{A}_5^2=60$ số tự nhiên thỏa mãn.
}
\end{ex}

\begin{ex}%[0H9V4-3]%[Dự án đề kiểm tra Toán 11 GHKII NH23-24- Tổng Nguyễn]%[THPT Đống Đa - Hà Nội]
Trong mặt phẳng $Oxy$, viết phương trình đường tròn $(C)$ có bán kính $R=5$ tiếp xúc với đường thẳng $d\colon 3x+4y+1=0$ tại điểm $A(1;-1)$.
% \shortans[]{$0{,}4$}
\loigiai{Đường thẳng $d$ có vectơ pháp tuyến $\overrightarrow{n}=(3;4)$.\\
Đường thẳng $\Delta$ đi qua $A(1;-1)$ và vuông góc với $d$ nên có vectơ chỉ phương $\overrightarrow{u}=(3;4)$. Phương trình của $\Delta$ là
\[\heva{&x=1+3t\\&y=-1+4t} \quad t \in \mathbb{R}.\]
Đường tròn $(C)$ có bán kính $R=5$,  tiếp xúc với $d$ tại $A(1;-1)$ nên tâm $I$ thuộc $\Delta$. Suy ra $I(1+3t; -1+4t)$.\\
Ta có \[ IA=R\Leftrightarrow \sqrt{(-3t)^2+(-4t)^2}=5\Leftrightarrow 25t^2=25 \Leftrightarrow  \hoac{&t=1\Rightarrow I_1(4;3) \ \text{(loại)}\\&t=-1\Rightarrow I_2(-2;-5)\ \text{(nhận)}.}\]
Vậy   $\dfrac{a}{b}=\dfrac{-2}{-5}=\dfrac{2}{5}=0{,}4$.
}
\end{ex}

\begin{ex}%[0D8V2-7]
Một giảng đường có $20$ hàng ghế, mỗi hàng có $24$ ghế. Trong buổi sinh hoạt công dân đầu khóa, có $444$ sinh viên tham dự và ngồi tại giảng đường này. Hỏi có ít nhất bao nhiêu hàng ghế có số sinh viên ngồi như nhau (nhập đáp án vào ô trống)?
% \shortans{5}
\loigiai{
Ta thấy rằng các hàng ghế có số sinh viên ngồi khác nhau đôi một thì số ghế trống ở các hàng ghế đó cũng đôi một khác nhau.\\
Tổng số ghế trong giảng đường là $20\cdot 24=480$ (ghế), suy ra số ghế trống là $480-444=36$ (ghế).\\
Gọi $k_0,\, k_1,\, \ldots,\, k_n$ lần lượt là số hàng ghế có $0,\, 1,\, \ldots,\, n$ chỗ trống.\\
Khi đó $0\cdot k_0 + k_1+2k_2+\cdots + nk_n =36$ và $k_0+k_1+\cdots +k_n=20$.\\
Các hàng ghế có số sinh viên ngồi khác nhau đôi một thì số ghế trống ở các hàng ghế đó cũng đôi một khác nhau.\\
Nếu số hàng có số ghế trống bằng nhau không vượt quá $4$ thì có ít nhất $5$ giá trị $k_i$ sao cho $1\leq k_i \leq 4$.\\
Mặt khác $4\cdot (0+1+2+3+4)=40>36$ nên trường hợp này không thể xảy ra.\\
Ta có $5\cdot (0+1+2)+4\cdot 3+9=36$. Do đó tồn tại $5$ hàng ghế mà có số ghế trống bằng nhau hay tồn tại $5$ hàng ghế có số sinh viên ngồi bằng nhau. Ta có thể phân bổ sinh viên vào các hàng ghế như sau:
Xếp $5$ hàng ghế mà mỗi hàng có $24$ sinh viên, $5$ hàng ghế mà mỗi hàng có $23$ sinh viên, $5$ hàng ghế có mà mỗi hàng $22$ sinh viên, $4$ hàng ghế mà mỗi hàng có $21$ sinh viên, $1$ hàng ghế có $15$ sinh viên.\\
Vậy số hàng ghế có số sinh viên ngồi như nhau ít nhất có thể đạt được là $5$ hàng.
}
\end{ex}
\Closesolutionfile{ans}
\Closesolutionfile{ansbook}
